\section{Methodology}
\label{sec:methodology}
Our proposed approach consists of three layers: static analysis, dynamic analysis, and agentic analysis. Outputs from the first two layers are fed into the final layer for semantic reasoning about the intent of a given skill. Figure~\ref{fig:system_architecture} illustrates the overall pipeline. The final layer outputs its findings in Static Analysis Results Interchange Format (SARIF) to enable seamless integration with other security tools. In this way, we provide not only triggered alerts but also a maliciousness score that summarizes the evidence and severity information from all analysis layers for each skill. An example SARIF output generated by our scanning approach is provided in Appendix~\ref{app:full_pipeline_sarif_example}.
% to ensure seamless integration with existing Static Application Security Testing (SAST) infrastructure, 

\begin{figure*}[!htbp]
    \centering
    \includegraphics[
        width=\textwidth
    ]{prism-uploads/system_architecture_short.pdf}
    \caption{Overview of the proposed framework.
Human review is an optional escalation step for cases
that remain unresolved after three reasoning rounds.
All 140 held-out test predictions were produced
automatically without human adjudication.}
    \label{fig:system_architecture}
\end{figure*}
% if the computed severity score equals or exceeds the 0.8 threshold. 
%The overall efficacy of the Agentic layer is systematically evaluated using standard classification metrics, specifically True Positives (TP), False Positives (FP), True Negatives (TN), False Negatives (FN), Accuracy, Precision, and the F1-score.

\subsection{Static Analysis Layer}
\label{sec:jiffylabs-static}
Our static layer performs a lightweight analysis of agent skills without executing them. Its purpose is to quickly identify potentially malicious patterns through syntactic and structural inspection before execution. To do so, we adopt a rule-based approach combining
regular-expression rules from Microsoft's
OSS-Detect-Backdoor~\cite{microsoft_ossgadget}
and YARA rule collections identified through the
cybersecurity-dev Awesome YARA resource
list~\cite{cybersecurity_dev_awesome_yara}.

OSS-Detect-Backdoor and YARA provide general-purpose
scanning capabilities rather than dedicated detection
of malicious agent skills. We use their rules to identify
potentially suspicious patterns in skill artifacts.
A rule match is treated as supporting evidence for
further analysis, rather than sufficient evidence
that a skill is malicious.

The initial inventory contained 115,209 rules across
the regular-expression and YARA collections combined.
We removed duplicate rules and rules identified as
unusable during curation, retaining a combined inventory
of 58,674 rules. Separately from rule removal, we reviewed
the severity assignments according to the security
relevance of the matched patterns to our threat model.
Rules considered less security-relevant were assigned
lower severity, while rules considered more
security-relevant were assigned higher severity.
Severity reassignment changes the interpretation of
retained findings rather than the number of retained rules.

The retained rules are assigned to three severity levels,
as shown in Table~\ref{tab:severity_levels}.
Rather than collapsing the findings into a single static
risk score, we preserve each finding's severity and
provenance as supporting evidence for the Agentic
Analysis Layer (Section~\ref{sec:jiffylabs-llm}),
which jointly reasons over static, dynamic, and
contextual evidence to produce the final verdict.

% To do so, we adopt a rule-based approach that includes regular-expression and YARA rules. The regular-expression rules are collected from Microsoft's OSS-Detect-Backdoor~\cite{microsoft_ossgadget}, while the curated YARA rules are sourced from publicly available repositories~\cite{yara}. 

% OSS-Detect-Backdoor and YARA provide general-purpose
% scanning capabilities rather than dedicated detection
% of malicious agent skills. We use their rules to identify
% potentially suspicious patterns in skill artifacts.
% A rule match is treated as supporting evidence for further
% analysis, rather than sufficient evidence that a skill
% is malicious. We therefore manually reviewed and filtered the initial rule set based on its applicability to our threat model. Specifically, we retained rules that target high-risk behaviors, apply to typical agent-skill environments, and minimize false positives in benign automation scripts. 

% Of the 115,209 initial rules, we retained 58,674 applicable rules. To provide a more nuanced risk assessment, we categorize each detection rule into one of three severity levels, as shown in Table~\ref{tab:severity_levels}. Rather than collapsing these findings into a single static risk score, we preserve each finding's severity and provenance as supporting evidence for the Agentic Analysis Layer (Section~\ref{sec:jiffylabs-llm}), which jointly reasons over static, dynamic, and contextual evidence to produce the final verdict.

\begin{table}[!t]
\centering
\caption{\textsc{Severity levels assigned to static analysis findings.}}
\small
\renewcommand{\arraystretch}{1.2}
\begin{tabularx}{\columnwidth}{>{\centering\arraybackslash}p{1.4cm}>{\RaggedRight\arraybackslash}X}
\toprule
\textbf{Severity} & \textbf{Description} \\
\midrule
Low &
Common programming patterns or indicators frequently found in benign software. Examples include access to standard system paths (e.g., \texttt{/usr/local/include}). These signals are not malicious on their own but provide contextual information. \\

\rowcolor[HTML]{F2F2F2}
Medium &
Dual-use behaviors that may be legitimate but are also commonly observed in malware, such as manual socket creation or unusual data encoding. \\

High &
% Assigned to strong indicators of malicious intent. These include behaviors rarely found in benign software, such as environment variable exfiltration or anti-analysis techniques. A single high-severity trigger significantly impacts the final assessment. \\
Strong indicators, including behaviors rarely found in benign software, such as environment variable exfiltration or anti-analysis techniques. A single high-severity trigger significantly affects the final assessment. \\
\bottomrule
\end{tabularx}
\label{tab:severity_levels}
\end{table}


% \begin{table}[!htb]
% \caption{Severity levels and corresponding weights used in the static analysis scoring system.}
% \centering
% \small
% \renewcommand{\arraystretch}{1.2}
% \begin{tabularx}{\textwidth}{ccX}
% \hline
% \textbf{Severity} & \textbf{Weight} & \textbf{Explanation} \\
% \hline
% Low & 40 &
% Assigned to common programming patterns or indicators frequently found in benign software. Examples include accessing standard system paths (e.g., \texttt{/usr/local/include}) or detecting generic HTTPS URLs. These signals are not malicious on their own but provide contextual information. \\
% \hline
% Medium & 70 &
% Assigned to suspicious or dual-use behaviors that may be legitimate but are also commonly observed in malware, such as manual socket creation or unusual data encoding. \\
% \hline
% High & 100 &
% Assigned to strong indicators of malicious intent. These include behaviors rarely found in benign software, such as environment variable exfiltration or anti-analysis techniques. A single high-severity trigger significantly impacts the final assessment. \\
% \hline
% \end{tabularx}
% \label{tab:severity_weights}
% \end{table}


%To quickly assess whether a skill is malicious, we introduce a \textit{Static Risk Score}, denoted as $S_{\mathrm{static}}$, which aggregates all triggered rules into a single value ranging from \num{0} to \num{100}. A higher score indicates a greater likelihood that the skill is malicious. To avoid an averaging effect in which critical high-risk signals are diluted by numerous low-severity matches, we adopt a weighted hybrid scoring scheme:


%\begin{equation}
% S_{\mathrm{static}} =
% 0.8 \cdot \max(\text{severity})
% + 0.2 \cdot \text{avg}(\text{severity})
% \label{eq:mal_score}
% \end{equation}

% The scoring formulation in Equation~\ref{eq:mal_score} prioritizes the most severe detected behavior, ensuring that a single high-risk indicator (e.g., a backdoor-related rule) significantly increases the overall score. At the same time, the average component captures the cumulative effect of multiple lower-severity signals, allowing skills with consistently suspicious patterns to be ranked higher than those with only isolated findings.

\subsection{Dynamic Analysis Layer}
\label{sec:jiffylabs-dynamic}

In this layer, we execute agent skills inside a controlled sandbox to generate runtime evidence for downstream dynamic analysis. The sandbox uses the prompt in Appendix~\ref{app:sandbox_configuration} to inspect \texttt{SKILL.md} and associated local files, create harmless dummy inputs when necessary, and exercise safe local workflows such as setup checks, validation commands, dry runs, mock executions, help or version commands, and tests. Rather than fully executing potentially harmful operations, the sandbox prevents actions that could access real credentials, exfiltrate data, install persistence, bypass security controls, execute unknown remote payloads, or perform real-world account operations. When such an action is encountered, it is intentionally skipped, and the requested command or instruction together with the reason for blocking it is recorded. This safety-constrained execution preserves evidence of malicious intent without requiring the harmful action itself to succeed. 

During sandbox execution, we collect four complementary categories of runtime evidence:
executed and blocked agent actions, system-call and process activity, network behavior,
and filesystem changes. 
We parse these artifacts to extract commands, actions reported
as skipped for safety reasons, security-related statements,
execution failures, process and file-access activity,
network observations, and filesystem changes. These extracted
observations form the behavioral features used by the
downstream attribution and scoring steps. Statements in the
execution agent's output are distinguished from actions
observed in system-call, network, or filesystem records:
a reported refusal or skipped command does not establish
that the corresponding behavior occurred. The execution
agent's output is therefore an evidence source, not the
final dynamic verdict.
% We process these artifacts and normalize them
% into behavioral features. The parsers extract executed commands, skipped
% unsafe actions, security assertions, and execution failures; reconstruct process and
% file-access activity; and identify network observations and filesystem changes. The
% sandbox output is treated as one source of evidence rather than as the final dynamic
% verdict.
Although downstream parsing and scoring follow fixed rules,
the execution agent determines which permitted actions to
attempt. The collected evidence may therefore vary across
executions. A single sandbox attempt does not establish
complete behavioral coverage, and the absence of an observed
malicious action is not sufficient to conclude that a skill
is benign.

We then examine the origin and context of each observation
to distinguish instructions contained in skill artifacts
from descriptions introduced during analysis. For example,
an execution agent's statement that it refused to transmit
credentials is not evidence that credential transmission
occurred. If the corresponding instruction can be traced
to \texttt{SKILL.md} or another skill artifact, however,
it may support an assessment of dangerous intent even
when the action was not executed.

Similarly, security terminology in defensive documentation,
test fixtures, or refusal messages is not treated as direct
evidence of malicious behavior merely because it mentions
an attack. Dependency failures indicate execution limitations
rather than maliciousness by themselves. This attribution
step is intended to reduce confusion between observed
actions, source-grounded instructions, and analysis-generated
descriptions.
% We then apply a provenance-aware evidence attribution technique to distinguish skill-attributable behavior from security terminology introduced by the analysis environment. For example, descriptions of a blocked attack, defensive documentation, test or validation fixtures, safe refusals, and dependency failures are not treated as direct malicious evidence. In contrast, instructions explicitly attributable to \texttt{SKILL.md} or another skill artifact can be promoted when they contain concrete dangerous primitives. This distinction prevents descriptive security language from being confused with malicious instructions.

Next, we combine related observations across evidence sources by identifying higher-level relationships, such as a remote download followed by executable preparation and process execution.
% , sensitive-file access combined with external transmission, or permission-bypass instructions associated with unsafe actions. 
Finally, we apply patterns developed on the development
split for skill-attributable prompt manipulation, permission bypass, download-and-execute behavior, untrusted software delivery, and file-upload or exfiltration activity.

Each detected pattern is assigned a base severity score of 40, 70, or 100 for low-,
medium-, or high-severity behavior, respectively. The base score is adjusted according
to the quality and provenance of the evidence, as summarized in
Table~\ref{tab:dynamic_evidence_weights}. We then compute an effective score for each pattern
$e_i \in [0,100]$ from its base severity score and evidence multiplier before applying the
corresponding guardrails. We additionally cap scores from noise and supporting pattern
families to prevent weak dual-use behavior from dominating the verdict, while directly
executed high-risk malicious patterns receive a score of 100. 


The dynamic score is a heuristic summary of the available
evidence, not a calibrated probability of maliciousness.
The scoring design gives the strongest pattern the largest
contribution so that a severe observation is not diluted
by weaker findings. Additional patterns contribute at a
reduced weight, while a bounded bonus accounts for support
from independent attributable evidence groups. Deduplication
limits repeated contributions from observations referring
to the same underlying assertion. The evidence multipliers
and guardrails distinguish weak narrative evidence,
source-grounded intent, and observed execution.

Formally, let $E=\{e_1,e_2,\ldots,e_n\}$ denote the effective scores of the unique correlated patterns and let $e_k=\max_i e_i$ be the dominant signal. The detector aggregates these scores as

% \begin{table}[!t]
% \centering
% \caption{Evidence types and scoring adjustments used by the dynamic detector.}
% \label{tab:dynamic_evidence_weights}
% \scriptsize
% \renewcommand{\arraystretch}{1.08}
% \setlength{\tabcolsep}{3.5pt}

% \begin{tabularx}{\textwidth}{@{}l c c X@{}}
\begin{table*}[!t]
\centering
\caption{\textsc{Evidence types and scoring adjustments used by the dynamic detector.}}
\label{tab:dynamic_evidence_weights}
\small
\renewcommand{\arraystretch}{1.08}
\setlength{\tabcolsep}{2pt}

\begin{tabularx}{\textwidth}{l c c >{\RaggedRight\arraybackslash}X}
\toprule
\textbf{Evidence Type}
& \textbf{Mult.}
& \textbf{Guardrail}
& \textbf{Interpretation} \\
\midrule

Ungrounded narrative
& 0.45
& $\leq 25$
& Security-related narrative without concrete source grounding. \\

\rowcolor[HTML]{F2F2F2}
Supporting intent only
& 0.50
& $\leq 25$
& Weak supporting evidence not observed at runtime. \\

Intent only
& 0.55
& $\leq 30$
& Unexecuted intent without sufficient source grounding or corroboration. \\

\rowcolor[HTML]{F2F2F2}
Grounded intent
& 0.70
& --
& Intent grounded in a concrete skill source, command, path, or code primitive. \\

Strong intent only
& 0.85
& $\geq 85$
& Source-grounded dangerous primitive identified but safely blocked. \\

\rowcolor[HTML]{F2F2F2}
Executed
& 1.00
& --
& Behavior directly observed during sandbox execution. \\

Corroborated
& 1.15
& --
& Executed behavior supported by multiple independent runtime artifacts. \\

\bottomrule
\end{tabularx}
\end{table*}




\begin{equation}
S_{\mathrm{raw}}
=
\min\left(
100,\,
e_k
+
\lambda\sum_{i\neq k}e_i
+
B(C)
\right),
\label{eq:dynamic_score}
\end{equation}

where $\lambda=0.1$ limits the contribution of additional correlated patterns, and $B(C)\in\{0,5,10,15\}$ is a saturated bonus determined by the number $C$ of independent attributable evidence groups. The bonus is disabled when directly executed malicious runtime behavior has already been observed. Evidence groups are deduplicated before counting to avoid inflating the score when several patterns originate from the same underlying assertion. The current evaluation does not include a sensitivity
analysis of these coefficients and guardrails; their
individual effects on detection performance therefore
remain unquantified.

% A final safeguard prevents collections consisting only of non-decisive narrative, supporting, or ungrounded intent evidence from crossing the malicious decision boundary. In contrast, verified source-grounded dangerous intent and directly executed malicious behavior remain high-confidence signals. The finalized detector classifies a skill as malicious when

% \begin{equation}
% S_{\mathrm{dyn}} \geq 35.
% \label{eq:dynamic_threshold}
% \end{equation}

% The dynamic rules and the threshold of 35 were finalized using only the 560-sample development split and were then frozen before evaluation. Thus, the LLM used during sandbox execution generates one source of runtime evidence, while the final dynamic verdict is produced by deterministic parsing, provenance attribution, evidence correlation, and scoring.


\subsection{Agentic Analysis Layer}
\label{sec:jiffylabs-llm}

% Answer the question: Why we need Agentic layer detection in our pipeline
While preceding analysis layers are highly specialized,
they exhibit inherent limitations. Our rule-based static
analysis effectively identifies syntactic anomalies but
has limited contextual awareness, potentially failing
to distinguish between benign administrative operations
and malicious activities (for instance, API-key access
and outbound requests used for legitimate authentication
or data exfiltration), thereby contributing to false
positives. Conversely, dynamic analysis yields concrete
behavioral traces but is constrained by execution-path
coverage and sandbox safety restrictions, which may
leave dormant threats or logic bombs undetected
within sandbox environments.

To address these issues, we implemented an Agentic layer that serves as the primary semantic-reasoning engine and
final decision-making component of the detection pipeline. It aggregates signals from
the static and dynamic layers and enriches the analyzed artifact with context retrieved
from threat-intelligence sources and a vector database. These sources include online
reports of malicious skills and established security guidance, such as OWASP
recommendations~\cite{owaspagentskills}. The retrieved context provides information on
emerging attack techniques, known vulnerabilities, and malicious domains. The agentic analysis operates as a compound AI system that uses Group Chat Orchestration~\cite{orchestration_design_pattern}.


% We leverage Large Language Model models to assess the true intent underlying the source code. 
%This semantic querying operates in the millisecond to second range, enabling the pipeline to expose obfuscated malicious intent with a precision comparable to behavioral analysis, while maintaining the high throughput necessary for real-time integration and large-scale scanning. Furthermore, integrating RAG significantly mitigates LLM hallucination. Standard LLMs are constrained by their training data cut-off dates; enveloping the LLM with our proprietary vulnerability reports ensures the system remains adaptive to emerging attack vectors within the software supply chain, particularly in the dynamic AI skills market.
%

% We introduce the Agentic layer to mitigate these bottlenecks by facilitating rapid, context-aware intent inference. 


% Answer the question: How RAG work
% Our semantic-reasoning layer is coordinated by a Planner Agent, which decomposes structural package investigations into discrete, actionable tasks and assigns them to specialized sub-agents. These sub-agents execute the tasks and report their findings to the planner. For example, suspicious lines identified by the static analysis layer serve as focal points for further investigation. The sub-agents operate according to the Reasoning and Acting (ReAct) framework~\cite{yao2023react}. Rather than following fixed workflows, each agent iteratively analyzes the current context, determines the next step, and invokes the appropriate analytical tools. This adaptability enables the agents to navigate obfuscated code or ambiguous structures based on intermediate tool outputs while continuously enriching a shared knowledge state.
Our semantic-reasoning layer is coordinated by a Planner
Agent and uses a predefined set of sub-agents with fixed
roles. The planner decomposes each skill investigation
into tasks and assigns them to the existing sub-agents.
The evaluated implementation reuses these predefined
agents rather than creating new agent roles for
individual skills, reflecting resource constraints.
Fixed roles do not imply a fixed sequence of investigative
actions. Within their assigned roles, the sub-agents
operate according to the Reasoning and Acting (ReAct)
framework~\cite{yao2023react}. Each sub-agent iteratively
examines the current context and intermediate tool
outputs to determine its next action.
For example, suspicious lines identified by the static
analysis layer can serve as focal points for a delegated
task. The sub-agents report their findings to the planner
for evidence assessment.

To improve architectural resilience, we implement fallback workflows for cases where
evidence from earlier analysis layers is unavailable or incomplete, such as when dynamic
sandbox execution times out or a static scan produces no findings (e.g., no suspicious
lines are flagged). In these cases, the Contextual RAG Agent autonomously retrieves and
interprets information from the available RAG context and the \texttt{SKILL.md} file,
allowing the analysis to continue when static or runtime evidence is unavailable.

% Next, an Evaluator Agent assesses the aggregated evidence
% and produces the final verdict using the semantic findings,
% static alerts, and dynamic evidence. Findings are interpreted
% in context: a security-sensitive operation does not by itself
% establish malicious intent, while an action reported as
% blocked is distinguished from an action observed during
% execution. Each verdict is required to cite supporting
% evidence from skill artifacts or execution records.
% These citations support inspection of the decision but do
% not guarantee that the interpretation is correct.
Next, the Evaluator Agent assesses the aggregated evidence,
including semantic findings, static alerts, and dynamic
evidence. Findings are interpreted in context:
a security-sensitive operation does not by itself
establish malicious intent, while an action reported
as blocked is distinguished from an action observed
during execution. Based on this assessment, the
Evaluator Agent weighs the available evidence,
determines the final verdict and risk score, and
generates a structured report. Each verdict must
cite supporting evidence from skill artifacts or
execution records so that the decision is traceable
and can be independently inspected. These citations
support transparency but do not guarantee that
the interpretation is correct. Appendix~\ref{app:agent-instructions} summarizes the
instructions for the Planner and Evaluator.

% Next, an Evaluator Agent assesses the aggregated evidence. It formulates the final verdict
% by combining the semantic analysis produced within this layer with findings and severity
% scores from the earlier scans. To reduce hallucinations, each verdict must be supported by
% verifiable evidence and source citations drawn directly from package artifacts and
% execution results.
The evaluation proceeds iteratively, with stopping decisions
based on the assessed strength and completeness of the
available evidence.
% Our approach adopts an adaptive evaluation strategy based on a dynamic confidence
% mechanism.
Early stopping is triggered when the system identifies severe, unambiguous
malicious behavior (e.g., explicit remote code execution or credential exfiltration)
with high confidence or when elevated dynamic severity scores provide sufficient
evidence of malware. 
The system then halts further investigative iterations,
and the Evaluator Agent finalizes the verdict,
risk score, and structured report.
Conversely, when contextual evidence is ambiguous, incomplete, or indicative of dual-use
functionality and the dynamic severity scores remain inconclusive, the system continues
the evaluation and returns to the Planner Agent to retrieve additional context.

To prevent unbounded execution, iterative reasoning is limited to three retrieval rounds. If the
system reaches this limit without a definitive verdict, it halts autonomous
classification and escalates the case for human review, allowing a security analyst to
examine the aggregated evidence and make the final determination.

Although the architecture supports escalation after three unresolved
reasoning rounds, no held-out test sample required human adjudication.
Therefore, all predictions reported in
Tables~\ref{tab:rag_evaluation_metrics} and~\ref{tab:full_vs_baselines}
were produced automatically by the pipeline without human modification
of the final label.

%Rather than relying on rigid termination points,


%The documents are indexed in Qdrant, a high-performance vector database. 
% To manage extensive threat reports, we use a semantic chunking strategy that partitions lengthy documents into smaller, semantically coherent segments. This approach preserves the core meaning while optimizing the context window for the research agents.
% Retrieval effectiveness is maintained through a combination of top-$k$ retrieval and hybrid search methods that leverage both dense vector embeddings and sparse keyword matching. 
%Advanced database reranking mechanisms are left for future work. Currently, the architecture operates without a strict similarity threshold because of the limited size of the reference dataset.
% Answer the question: Why we choose RAG instead of other approaches.
% Compared to alternative methodologies such as LLM fine-tuning or long-context prompting, our RAG-based approach demonstrates distinct structural advantages. Table \ref{tab:llm_methods_comparison} provides a comprehensive comparison of these paradigms.

% \begin{table}[ht]
% \centering
% \caption{Comparative Analysis of LLM Integration Methodologies}
% \label{tab:llm_methods_comparison}
% \small
% \renewcommand{\arraystretch}{1.3}
% \begin{tabularx}{\textwidth}{@{}>{\raggedright\arraybackslash}p{0.18\textwidth} Y Y Y@{}}
% \toprule
% \textbf{Evaluation Metric} & \textbf{LLM Fine-tuning} & \textbf{Long-Context Prompting} & \textbf{RAG (Our Approach)} \\ \midrule
% \textbf{Knowledge Updatability} & Static (Requires complete model retraining for new threats) & Dynamic (Context must be provided at runtime) & Highly Dynamic (Real-time updates via vector database) \\
% \textbf{Resource Overhead} & Very High (Requires significant training compute and dataset curation) & High (Inference costs scale linearly or quadratically with token length) & Moderate (Efficient vector search combined with bounded inference) \\
% \textbf{Hallucination Mitigation}& Low (Answers remain bounded by training data distribution) & Moderate (Vulnerable to "lost in the middle" information degradation) & High (Responses are strictly grounded in retrieved threat intelligence) \\
% \textbf{Scalability} & Low (Impractical for continuous, daily vulnerability updates) & Moderate (Constrained by the model's maximum context window) & High (Seamlessly scales with the underlying Qdrant database) \\
% \bottomrule
% \end{tabularx}
% \end{table}

% Answer the question: What data source are retrived

% Answer the question: How are documents and index are stored


% Answer the question: How retrieval quality ensured



% Answer question: What model this LLM is using, why we use this model
%We deploy Claude 4.6 Sonnet~\cite{claude_sonnet_46} for complex, reasoning-intensive tasks and Claude 4.5 Haiku~\cite{claude_models_overview} for simpler, routine operations. We use this model combination for malicious skill detection to balance reasoning capability with computational efficiency. This dual-model approach optimizes both cost and performance. In particular, Claude 4.6 Sonnet provides strong logical-deduction and advanced code-comprehension capabilities, which are essential for accurately analyzing configuration files, such as \texttt{SKILL.md}, and identifying obfuscated or sophisticated malicious intent by correlating aggregated static and dynamic evidence from preceding layers. In contrast, delegating routine tasks, such as preliminary data extraction, formatting, and basic intent checking, to Claude 4.5 Haiku significantly reduces latency and token expenditure.

% Supply-chain analysis requires processing extensive datasets across multiple iterative loops. Delegating routine tasks, such as preliminary data extraction, formatting, and basic intent checking, to Claude 4.5 Haiku significantly reduces latency and token expenditure. Meanwhile, Claude 4.6 Sonnet is reserved for critical, high-level reasoning roles, including the Planner, Evaluator, and Synthesizer agents. Claude 4.6 Sonnet provides strong logical deduction and advanced code-comprehension capabilities, which are essential for accurately analyzing configuration files, such as \texttt{SKILL.md}, and identifying obfuscated or sophisticated malicious intent by correlating aggregated static and dynamic evidence from preceding layers. This tiered model strategy aligns with our compound AI architecture and the Group Chat Orchestration Pattern, enabling dynamic subtask routing so that computational resources are allocated according to the cognitive demands of each task and are not wasted on low-complexity operations.

%This dual-model approach optimizes cost and performance. Supply chain analysis inherently requires processing extensive datasets across multiple iterative loops. Delegating routine tasks—such as preliminary data extraction, formatting, and basic intent checking—to Claude 3.5 Haiku significantly reduces latency and token expenditure. Meanwhile, Claude 4.6 Sonnet is strictly reserved for critical, high-level cognitive roles, including the Planner, Evaluator, and Synthesizer agents.

%We utilize this specific model combination for the malicious skill detection problem for the following reasons:



% Answer the question: How the Agentic layer interact with LLM


%The LLM-based layer serves as the final semantic reasoning stage of the detection pipeline. While the static layer identifies suspicious syntactic patterns and the dynamic layer observes runtime behavior, the LLM layer evaluates these findings in context. Its goal is to distinguish truly malicious intent from benign documentation, security examples, setup scripts, or dual-use automation patterns.

%This layer operates on the original skill package together with evidence produced by the previous stages, including static rule matches, extracted suspicious lines, sandbox execution summaries, skipped unsafe commands, filesystem events, and network indicators. When local evidence is insufficient, the system can retrieve additional context from internal databases or external sources, such as reputation information for embedded URLs, package names, domains, or referenced tools. This layer consists of five main steps: Package ingestion and evidence preservation, Evidence aggregation, Context-aware LLM assessment, Context retrieval loop, and Final scoring and escalation. This design allows the pipeline to combine rule-based precision with semantic interpretation. In particular, the LLM layer helps reduce false positives from keyword-based detections while still preserving conservative handling of ambiguous or high-impact behaviors.

%As the first step, given an agent skill, we store its raw content. This preserves the original artifact for later auditing and allows the LLM to inspect the full context rather than only isolated alerts. Then, we run the static and dynamic analysis modules to collect the analysis findings and consolidate them into a structured evidence bundle. This bundle includes triggered rules, severity scores, relevant code or instruction snippets, observed runtime behavior, skipped dangerous actions, and any supporting contextual information retrieved from databases or web sources. Then, we run an LLM model to evaluate the aggregated evidence through a set of decision criteria. It checks whether the skill demonstrates clear malicious behavior, whether suspicious indicators are better explained by benign documentation or normal development workflows, and whether the available context is sufficient to make a confident judgment. If the evidence is ambiguous, the system retrieves additional context before producing a final verdict. This may include extracting more surrounding text from the package, checking the reputation of URLs or domains, examining referenced dependencies, or reviewing related metadata. The enriched context is then passed back to the LLM for re-evaluation. Finally, if the LLM identifies strong malicious intent or behavior, the package receives a high-risk assessment. If the evidence is clearly benign, it receives a low-risk assessment. If the system reaches the retry limit and the evidence remains inconclusive, the package is assigned a medium-risk score and flagged for manual review.

